後十九日復上書

韓愈

二月十六日，前鄉貢進士韓愈，謹再拜言相公閣下。

向上書及所著文後，待命凡十有九日，不得命。恐懼不敢逃遁，不知所為。乃復敢自納於不測之誅，以求畢其說，而請命於左右。

愈聞之，蹈水火者之求免於人也，不惟其父兄子弟之慈愛，然後呼而望之也。將有介於其側者，雖其所憎怨，茍不至乎欲其死者，則將大其聲疾呼而望其仁之也。彼介於其側者，聞其聲而見其事，不惟其父兄子弟之慈愛，然後往而全之也。雖有所憎怨，茍不至乎欲其死者，則將狂奔盡氣，濡手足，焦毛發，救之而不辭也。若是者何哉？其勢誠急，而其情誠可悲也。愈之強學力行有年矣。愚不惟道之險夷，行且不息，以蹈於窮餓之水火，其既危且亟矣，大其聲而疾呼矣，閣下其亦聞而見之矣。其將往而全之歟？抑將安而不救歟？有來言於閣下者曰：“有觀溺於水而爇於火者，有可救之道，而終莫之救也。”閣下且以為仁人乎哉？不然，若愈者，亦君子之所宜動心者也。

或謂愈，子言則然矣，宰相則知子矣，如時不可何？愈竊謂之不知言者，誠其材能不足當吾賢相之舉耳。若所謂時者，固在上位者之為耳，非天之所為也。前五六年時，宰相薦聞，尚有自布衣蒙抽擢者，與今豈異時哉？且今節度、觀察使，及防禦、營田諸小使等，尚得自舉判官，無間於已仕未仕者，況在宰相，吾君所尊敬者，而曰不可乎？

古之進人者，或取於盜，或舉於管庫。今布衣雖賤，猶足以方於此。情隘辭蹙，不知所裁，亦惟少垂憐焉。愈再拜。